\documentclass[11pt]{article}
\newcommand{\routelabel}{Expert}
\usepackage{eastbridge-handout}
\hypersetup{pdftitle={RPS Workshop: Expert}}
\begin{document}
\thispagestyle{plain}

\workshoptitle{Can you beat random play?}{Randomness and prediction}

\begin{context}
Against a bot that always chooses rock, you can win every throw by choosing paper. Against one that chooses each move independently with probability one third, any move you pick wins, loses and draws with equal probability. Its previous moves give you no help with the next one.

\end{context}

We'll calculate payoff as \textbf{+1 for a win, 0 for a draw, -1 for a loss}. Adding those values over a match gives your wins minus your losses. Whoever wins more throws wins the match.


\smallskip\begingroup\small
\setlength{\tabcolsep}{5pt}
\renewcommand{\arraystretch}{1.17}
\begin{tabularx}{\linewidth}{@{}YYYY@{}}
\toprule
\textbf{Your move} & \textbf{Opponent rock} & \textbf{Opponent paper} & \textbf{Opponent scissors} \\
\midrule
Rock & 0 & -1 & +1 \\
Paper & +1 & 0 & -1 \\
Scissors & -1 & +1 & 0 \\
\bottomrule
\end{tabularx}\endgroup\par\smallskip

\subsection*{The Nash equilibrium}

When both players choose independently and uniformly, neither can improve their expected payoff by changing strategy alone. This is a \textbf{Nash equilibrium}. Even a bot with the full match history and a very good prediction algorithm has expected net payoff zero against an independent uniform opponent.


Several house bots do something more predictable. One favors rock; another repeats a sequence; Copycat copies your last move. You can use those habits to predict what comes next and choose a winning reply. Much of the workshop is about learning to recognize them from a short history.


A repeating \code{rock, paper, scissors} bot is an especially helpful example. Count its moves after any complete cycle and you'll find equal totals. But once you've seen rock, you know paper is coming. A random bot can also produce that sequence by chance, so you'll need more than one occurrence before treating it as a pattern.


The kit uses seeded pseudorandom generators to make practice matches repeatable. The exercises use those seeds to compare versions under the same conditions. They assume each bot's current random choice is private; reconstructing an opponent's generator from its seed would be a different problem.
\par

\clearpage
\workshoptitle{What your bot sees}{Inside bot.py}

\begin{context}
For each match, the arena starts a fresh Python process and imports bot.py. It calls \code{setup(config)} once, if you've defined it, then asks \code{next\_move(...)} for each throw. Here is the starter bot:

\end{context}

\begin{lstlisting}[style=python]
from random import Random
from rpsdk import Move

_rng = Random()

def setup(config):
    _rng.seed(config["seed"])

def next_move(my_history, opponent_history, match_state):
    return _rng.choice(list(Move))
\end{lstlisting}

Return \code{Move.ROCK}, \code{Move.PAPER} or \code{Move.SCISSORS} from next\_move. These values also have a comparison method: \code{a.beats(b)} tells you whether a beats b. The arena accepts lowercase strings too, though the named values help catch spelling mistakes.


\subsection*{Reading the histories}

Both history lists contain \textbf{completed throws, oldest first}. On the first call, both are empty. After two throws, you might see:


\begin{lstlisting}[style=python,numbers=none,xleftmargin=4pt,framexleftmargin=0pt]
my_history       == [Move.ROCK, Move.PAPER]
opponent_history == [Move.SCISSORS, Move.ROCK]
\end{lstlisting}

You won both of those throws. \code{opponent\_history[-1]} gives their last move, while \code{my\_history[-1]} gives yours. Check that a list has an entry before using \code{[-1]}, or two entries before using \code{[-2]}. The histories never include the throw you're about to choose.


\smallskip\begingroup\small
\setlength{\tabcolsep}{5pt}
\renewcommand{\arraystretch}{1.17}
\begin{tabularx}{\linewidth}{@{}L{.28\linewidth}Y@{}}
\toprule
\textbf{State entry} & \textbf{Meaning} \\
\midrule
\code{round} & Zero-based throw index. The next call above has round 2. \\
\code{best\_of} & Total throws in this match, including draws. \\
\code{last\_outcome} & Your preceding throw's win, loss or draw; None initially. \\
\code{seed} & Your bot's reproducibility seed. \\
\code{timeouts}, \code{opponent\_timeouts} & Accumulated server timeouts in this match. \\
\bottomrule
\end{tabularx}\endgroup\par\smallskip

Each match starts with empty histories and fresh learned state. Reset any global counters in setup so your bot also works when tests call setup repeatedly. Seed the random generator there once; next\_move can then draw from it throughout the match.


\begin{checkpoint}
\textbf{Try it on paper:} Copycat copies your preceding move. Given the two histories above, what will it play next? Which history did you use?
\end{checkpoint}

\clearpage
\workshoptitle{Combine several strategies}{Route D / Expert}

\begin{context}
A frequency response, a Markov model and a Copycat response each work against different opponents. Suppose you keep all three running and adjust their weights as the match progresses. This route develops that idea using online learning, then looks at how the score near the end of a match can change your choice of move.

\end{context}

\lessonstep{Experiment 1}{an expert portfolio}

Build a small set of policies, called \textbf{experts}, that each return a probability distribution over rock, paper and scissors. Include uniform play, frequency and short-window responses, order-one and order-two models, and a reactive rule or two. A deterministic expert puts probability one on its chosen move.


Before throw $t$, expert $i$ proposes $p_{i,t}$. Give it weight $w_{i,t}$, normalize the weights, and sample your actual move from their weighted mixture. Store every proposed distribution before the opposing move is revealed.


After observing opposing move $b_t$, you can compute the one-step virtual gain of \textbf{every} proposal. Sum over the three possible actions $a$:


\[
g_{i,t} = \sum_a p_{i,t}(a)\,\operatorname{payoff}(a,b_t).
\]

RPS gives you \textbf{full-information feedback}: after seeing the opposing move, you can score every saved proposal, including those you didn't play. Score them first, then update their models for the next throw.


An exponential update is $w_{i,t+1}=w_{i,t}\exp(\eta g_{i,t})$. Start with equal positive weights. Use log weights and subtract the maximum before exponentiating to keep normalization stable. Every gain lies in $[-1,1]$.


\begin{checkpoint}
\textbf{First test:} use three fixed-action experts and one uniform expert. Feed them a short sequence you've written down and calculate the first weight update by hand. After normalization, check that the weights are finite, nonnegative and sum to one.
\end{checkpoint}

The standard library is enough for this implementation. Start with a few experts whose behavior you can follow, and add more once their updates work.


\textbf{Reading:} Freund and Schapire, \emph{A Decision-Theoretic Generalization of On-Line Learning and an Application to Boosting} (1997), DOI 10.1006/jcss.1997.1504. Arora, Hazan and Kale's 2012 survey develops the broader multiplicative-weights framework; its link is on the last page.
\par
\clearpage
\section*{A regret bound}

\begin{context}
For $K$ experts with gains in $[-1,1]$, the undiscounted exponential update gives the following bound. It holds on each realized history and compares the mixture's weighted gains with those of the best expert on that history:

\end{context}

\[
\max_i\sum_{t=1}^{T}g_{i,t}
  - \sum_{t=1}^{T}\sum_{i=1}^{K}\alpha_{i,t}\,g_{i,t}
\;\leq\; \frac{\log K}{\eta}+\frac{\eta T}{2}.
\]

Here $\alpha_{i,t}$ is the normalized weight. Setting $\eta=\sqrt{2\log K\,/\,T}$ makes the right side $\sqrt{2T\log K}$. If the opponent cannot see your current random draw, your conditional expected gain equals the weighted virtual gain. Your sampled moves will fluctuate around that expectation.


\lessonstep{Experiment 2}{include a uniform expert}

A uniform expert has virtual gain zero against every opposing move. Including it ensures that the best expert's cumulative virtual gain is at least zero. With $K=8$ and $T=501$, the bound allows a shortfall of about $0.091$ per throw. This comparison applies in expectation to the sampled policy; individual matches can finish below it.


Derive the bound by comparing the final log sum of weights with the weight of one expert. You'll need to bound each log moment-generating function for a variable in $[-1,1]$. Where does that range enter the constant?


\subsection*{Playing each expert separately}

An expert's virtual gains are calculated on the history your portfolio produced. If that expert had played the entire match, it could have produced a different history and caused an adaptive opponent to respond differently. The bound compares experts on the shared, realized history.


Try this against Cycle Counter, Markov and Adaptive Ensemble. Record the portfolio's result, then let each expert play fresh matches alone. How do those scores compare with its virtual gains inside the portfolio?


\begin{checkpoint}
\textbf{Derivation check:} identify the pathwise inequality and the expectation statement in your argument. What additional argument would you need to bound the probability of a large loss from sampling?
\end{checkpoint}

\textbf{Optional extension:} discount old gains and compare the result against Shifting Bias. The derivation above uses undiscounted sums; work out which steps would need to change for the discounted version.
\par
\clearpage
\section*{Mixing with uniform play}

\begin{context}
You can limit how much a learned policy affects your moves by mixing it with uniform play. Let $u$ be uniform and $r$ be your learned action distribution, and choose from:

\end{context}

\[
p=(1-\rho)u+\rho r, \qquad 0\leq\rho\leq1.
\]

Against an opponent that cannot see your current random draw, uniform play has conditional expected payoff zero and $r$ has payoff at least $-1$. The mixture therefore has conditional expected payoff at least $-\rho$ on each throw.


\lessonstep{Experiment 3}{vary the mixture}

Compare $\rho$ values $0$, $0.25$, $0.5$ and $1$ while holding the learned policy fixed. Smaller values give up some advantage when your model is right and reduce expected losses when it is wrong.


You could also choose $\rho$ using the sample count in the current context or the scores of stored predictions. Test that choice against an adaptive opponent. If you attach a statistical confidence bound, check that its assumptions allow dependence between observations.


\smallskip\begingroup\small
\setlength{\tabcolsep}{5pt}
\renewcommand{\arraystretch}{1.17}
\begin{tabularx}{\linewidth}{@{}L{.25\linewidth}YYYY@{}}
\toprule
\colhead{Variant} & \colhead{Stationary bias} & \colhead{Changing bias} & \colhead{Adaptive learner} & \colhead{Uniform} \\
\midrule
\rule{0pt}{22pt}Uniform &  &  &  &  \\
\rule{0pt}{22pt}Portfolio &  &  &  &  \\
\rule{0pt}{22pt}Portfolio mixed with uniform &  &  &  &  \\
\rule{0pt}{22pt}Discounted portfolio &  &  &  &  \\
\bottomrule
\end{tabularx}\endgroup\par\smallskip

Use fresh seeds for this comparison. Keep model state separate from the sampling RNG and reset both in setup. The logged proposal used for scoring should be exactly the one saved before the opposing move arrived.


\begin{summit}
\textbf{Challenge:} write an opponent that changes behavior in response to your past moves. Compare virtual expert gains with complete matches played by each expert. Can you produce a large difference between the two?
\end{summit}

So far we've optimized payoff per throw. The arena awards a match win for finishing ahead, which gives us another objective to consider.
\par
\clearpage
\section*{Choosing a move near the end}

\begin{context}
Suppose the opponent draws independently from a known, fixed distribution $q$. With one throw left and a one-point lead, a draw is enough to win the match. If you're one point behind, you need a win just to tie. We'll calculate how the score changes the best reply.

\end{context}

\lessonstep{Experiment 4}{a finite-horizon controller}

Let $V(n,d)$ be the best expected match points with $n$ throws remaining and current score difference $d$. At the end:


\[
V(0,d)=\begin{cases}
1 & d>0,\\
\tfrac12 & d=0,\\
0 & d<0.
\end{cases}
\]

For a candidate move $a$, $q$ determines probabilities $W(a)$, $D(a)$ and $L(a)$ of a throw win, draw and loss. The backward recursion is:


\[
\begin{aligned}
V(n,d)=\max_a\bigl[\,&W(a)\,V(n-1,d+1)\\
                  &+D(a)\,V(n-1,d)\\
                  &+L(a)\,V(n-1,d-1)\,\bigr].
\end{aligned}
\]

For a fixed known $q$, this expression is linear in the probabilities of your moves, so a pure action attains the maximum. Ties allow mixtures too. Modeling uncertainty about $q$ or an opponent that adapts to you would require extending the state and transition model.


Implement this for the final \textbf{30 throws}, using a precomputed table or memoized states. A full 501-throw dynamic program on every call would do much more work. Calculate $d$ as your wins minus losses; the remaining-throw count includes the move you're choosing now.


\subsection*{One throw left}

Take $q$ = 30\% rock, 60\% paper, 10\% scissors. The net-payoff maximizer is scissors ($+0.30$); paper earns $+0.20$. With one throw left and $d=0$, scissors also has the largest expected terminal points: $0.60+0.5\cdot0.10=0.65$. But with $d=1$:


\begin{itemize}
\item Paper loses with probability $0.10$, so its expected match points are $1-0.5\cdot0.10=0.95$.
\item Scissors loses with probability 0.30, so its expected match points are 0.85.
\end{itemize}

With that one-point lead, paper earns more expected match points. Calculate the result for rock as well, then repeat all three calculations with $d=-1$.


\begin{checkpoint}
\textbf{Test the controller:} check its terminal values and one-step choices against your calculations. Then play it against a fixed synthetic distribution and compare with a bot that maximizes payoff per throw. Once that works, try estimating $q$ from the match history.
\end{checkpoint}
\clearpage
\section*{Choose an experiment}

\begin{context}
Choose one of the questions below, or use a result from your earlier matches that you'd like to investigate. Leave time to run several seeds and save the version you tested.

\end{context}

\subsection*{Three possible investigations}

\textbf{Hidden order.} Compare context lengths one, two and three against Double Take, then add noise or a regime switch to your local test sequence. Measure the sample cost of discovering the longer dependency and the cost of retaining it after the regime changes.


\textbf{Adaptive opposition.} Compare an undiscounted portfolio, a discounted portfolio and its individual experts against Markov and Adaptive Ensemble. Explain the gap between one-step virtual gains and separate complete-match outcomes.


\textbf{Terminal utility.} Compare a final-30-throw controller with a net-payoff response. First use a known stationary $q$, then use an estimated $q$. Separate the effect of the objective from the effect of estimation error.


\subsection*{Report enough to reproduce it}

\begin{lstlisting}[style=shell]
rps-cli play --against double_take,markov,adaptive_ensemble \
  --best-of 501 --games 10 --seed 2000 --output evaluation.json
rps-cli package --output reviewed-bot.zip
\end{lstlisting}

Save the kit version, bot source, opponents, series length and seed range with the results. Report match outcomes and mean payoff per throw, using independent matches as your experimental units. If any runs have errors, fix those and rerun the comparison.


Before submitting, check weight normalization, empty contexts, repeated setup, tie-breaking and prediction timing. Run the bot on a lab machine too; its startup and move calculations need to fit the arena's time limits.


\begin{checkpoint}
\textbf{Finish:} submit your tested version and confirm that it's active. Show another participant the comparison you ran. Include the settings and source so they can try it themselves.
\end{checkpoint}

\subsection*{Further reading}

\begin{itemize}
\item Arora, Hazan and Kale (2012), \emph{The Multiplicative Weights Update Method: A Meta-Algorithm and Applications}. Theory of Computing 8, 121-164. \url{https://theoryofcomputing.org/articles/v008a006/}
\item Freund and Schapire (1997), \emph{A Decision-Theoretic Generalization of On-Line Learning and an Application to Boosting}. Journal of Computer and System Sciences 55, 119-139. \url{https://doi.org/10.1006/jcss.1997.1504}
\item Yale Open Courses, ECON 159, Lecture 8, mixed strategies and rock-paper-scissors. \url{https://oyc.yale.edu/economics/econ-159/lecture-8}
\end{itemize}

The first two references develop the online-learning results used in this route. The Yale lecture covers mixed strategies and the RPS equilibrium introduced in the shared pages.
\par
\end{document}
